As a motivation and running example we will consider the following objects:

FIG: Graph, VectorSp and Matrix.

Pretty soon we'll start thinking about cycles on graphs, dependent sets, bases, spanning sets, and in matrices we can think about columns. The concept of independence is not unique to graphs or vector spaces. These concepts were introduced by Whitney and Nakasawa back in 1930's. An itinerary for this will be 
\begin{enumerate}
    \item Basic intuitions and definitions.
    \item Non-representable matroids and operations.
    \item Positroids and characterizations.
    \item Matroids, posets and polytopes.
\end{enumerate}

\section{Independence}

\begin{Ex}
    For \(u,v,w=\dots\) Can be find \(a,b,c\in\bR\) such \(au+bv+cw=0\)? 
\end{Ex}

\begin{Def}
    A set of vectors is linearly independent if\dots Otherwise they're \term{dependent}.
\end{Def}

For graphs, consider the running example. A closed path that is minimally closed is a cycle. 
